\begin{tabular}{lccccc}
  \toprule
  \textbf{Variable} & \textbf{Region} & \textbf{Benchmark} & \textbf{Std ($\sigma$)} & \textbf{Lower Bound} & \textbf{Upper Bound} \\ 
  \midrule
  GDP & USA & 2.2094\% & --- & --- & --- \\ 
   & EZ & 1.3517\% & --- & --- & --- \\ 
   & CHN & 8.1924\% & --- & --- & --- \\ 
  CPI & USA & 2.0000\% & 1.7634\% & 0.2366\% & 3.7634\% \\ 
   & EZ & 2.0000\% & 1.8242\% & 0.1758\% & 3.8242\% \\ 
   & CHN & 3.0000\% & 1.9360\% & 1.0640\% & 4.9360\% \\ 
  Unemployment & USA & 5.6950\% & 1.9614\% & 3.7336\% & 7.6564\% \\ 
   & EZ & 9.0270\% & 1.5876\% & 7.4394\% & 10.6146\% \\ 
   & CHN & 4.4876\% & 0.3417\% & 4.1459\% & 4.8294\% \\ 
  \bottomrule
\end{tabular}
\par\vspace{4pt}
\begin{minipage}{\linewidth}
  \footnotesize \textit{Note}: Benchmark denotes the historical mean (GDP, Unemployment) or the official target (CPI) over the period 2000 to 2024. GDP thresholds are one-sided (stress if below benchmark); CPI and Unemployment thresholds are two-sided (benchmark $\pm$ one standard deviation).
\end{minipage}